\section*{अध्याय \ref{ch:b1:species-interactions} --- जातियों के बीच अन्योन्यक्रियाएँ}

\begin{solution}{exo:b1:species-interactions:1}
\omterm{def:b1:species-interactions:types}{स्पर्धा} $-/-$ (प्रकाश के लिए दो बलूत); \omterm{def:b1:species-interactions:types}{परभक्षण} $+/-$ (लिंक्स, खरगोश);
\omterm{def:b1:species-interactions:types}{परजीविता} $+/-$ (किलनी, हिरन); \omterm{def:b1:species-interactions:types}{सहोपकारिता} $+/+$ (मधुमक्खी, फूल);
\omterm{def:b1:species-interactions:types}{सहभोजिता} $+/0$ (छाल पर फ़र्न); अपकारिता $-/0$ (किसी बलूत की छाया के नीचे
घास)।
\end{solution}

\begin{solution}{exo:b1:species-interactions:2}
एक ही सीमाकारी संसाधन को एक ही ढंग से बरतती दो जातियाँ किसी स्थिर
परिवेश में अनिश्चित काल तक साथ नहीं रह सकतीं। पाँचों फुदकियाँ
वही भोजन बरतती हैं पर पेड़ के अलग-अलग भागों में और अलग-अलग
समयों पर: उनके \omterm{def:b1:ecosystem-organization:niche}{निकेत} भिन्न हैं, इसलिए उनमें से कोई दो ``एक ही ढंग'' नहीं हैं।
\end{solution}

\begin{solution}{exo:b1:species-interactions:3}
\omterm{def:b1:species-interactions:functional}{निपटान काल}: एक शिकार को पकड़ने, खाने और पचाने का समय, जिसके
दौरान कोई दूसरा नहीं लिया जाता। आक्रमण दर: प्रति इकाई समय खोजा गया
क्षेत्रफल गुणा पकड़ की प्रायिकता। अधिकतम दर $1/h = 2$ शिकार प्रति
घंटा।
\end{solution}

\begin{solution}{exo:b1:species-interactions:4}
ऐसी जाति जिसके हटने से \omterm{def:b1:ecosystem-organization:ecosystem}{समुदाय} की संरचना उसकी बहुलता के अनुपात से बाहर
बदल जाए, आम तौर पर किसी हावी स्पर्धी को काबू में रखने के कारण।
बहुलता उपस्थिति नापती है, असर नहीं; कोई दुर्लभ परभक्षी
जो होने वाले विजेता को खाता है उसका असर उसकी संख्या नहीं दिखाती।
\end{solution}

\begin{solution}{exo:b1:species-interactions:5}
\omterm{prop:b1:species-interactions:lotka}{सम-रेखा} 1: $N_1 + 0.5 N_2 = 500$ ($N_1$ पर अंतःखंड 500, $N_2$ पर
1000); \omterm{prop:b1:species-interactions:lotka}{सम-रेखा} 2: $N_2 + 0.6 N_1 = 400$ ($N_2$ पर 400, $N_1$ पर 667)।
वे इस तरह कटती हैं कि हर जाति स्वयं से अधिक सीमित हो ($\alpha = 0.5 <
K_1/K_2 = 1.25$; $\beta = 0.6 < K_2/K_1 = 0.8$): यानी
$N_1 = 429$, $N_2 = 143$ पर स्थायी सहअस्तित्व। $\alpha = 2$ के साथ:
\omterm{prop:b1:species-interactions:lotka}{सम-रेखा} 1 के अंतःखंड 500 और 250, दोनों \omterm{prop:b1:species-interactions:lotka}{सम-रेखा} 2 के (667 और 400) भीतर:
जाति 2 हमेशा जीतती है।
\end{solution}

\begin{solution}{exo:b1:species-interactions:6}
$f = aN/(1 + ahN)$: 10 पर, प्रति घंटा $2/(1 + 0.5) = 1.33$; 100 पर,
प्रति घंटा $20/(1 + 5) = 3.33$ (अधिकतम $1/h = 4$)। अर्ध-संतृप्ति
प्रति वर्ग मीटर $N = 1/(ah) = 20$ पर।
\end{solution}

\begin{solution}{exo:b1:species-interactions:7}
प्रकार III अनुक्रिया के साथ खाए गए शिकार का अंश नीचे घनत्व पर
घनत्व के साथ चढ़ता है: दुर्लभ शिकार अनुपात में कम लिए जाते हैं और
उबर सकते हैं। प्रकार II के साथ खाया गया अंश नीचे घनत्व पर सबसे ऊँचा
होता है (वहाँ परभक्षी कभी संतृप्त नहीं होता), इसलिए कोई छोटी
शिकार-आबादी और नीचे धकेल दी जाती है।
\end{solution}

\begin{solution}{exo:b1:species-interactions:8}
\emph{P.~caudatum} उसी जल-स्तंभ में वही बैक्टीरिया बरतता था
जो \emph{P.~aurelia}, और तेज़, अधिक किफ़ायती \emph{P.~aurelia}
ने भोजन को उससे नीचे ले जाया जितना उसे चाहिए था। \emph{P.~bursaria}
तली में भिन्न संसाधनों पर खाता था: \omterm{def:b1:ecosystem-organization:niche}{निकेत} ने तय किया, वृद्धि दर ने नहीं
--- \emph{P.~bursaria} \emph{P.~caudatum} से तेज़ नहीं बढ़ता।
\end{solution}

\begin{solution}{exo:b1:species-interactions:9}
परभक्षी उस ढर्रे से बिना हानि के अधिक बार मिलते हैं, उसे कम अच्छी तरह
सीखते हैं, और उस पर अधिक हमला करते हैं: अनुहारी और प्रतिरूप दोनों की
रक्षा गिर जाती है। इसलिए आम होते जाने के साथ अनुहारी की योग्यता गिरती है,
यानी एक ऋणात्मक बारंबारता-निर्भरता जो उसे प्रतिरूप के सापेक्ष दुर्लभ रखती है।
\end{solution}

\begin{solution}{exo:b1:species-interactions:10}
परभक्षी के बिना सीपी की \omterm{prop:b1:species-interactions:lotka}{सम-रेखा} हर स्पर्धी की \omterm{prop:b1:species-interactions:lotka}{सम-रेखा} से बाहर पड़ती
है: वह चट्टान जीत लेती है। \omterm{def:b1:species-interactions:types}{परभक्षण} सीपी का प्रभावी $K$ गिरा देता है
(उसकी \omterm{prop:b1:species-interactions:lotka}{सम-रेखा} भीतर की ओर खिंच जाती है) जब तक वह
दूसरों की \omterm{prop:b1:species-interactions:lotka}{सम-रेखाओं} को भीतर से न काट ले: तब स्पर्धी सीपी से अधिक
स्वयं से सीमित हो जाते हैं, और उसके साथ रहने लगते हैं। यह
परभक्षी सबसे मज़बूत स्पर्धी पर काम करता है और बाकियों के लिए जगह बनाता है।
\end{solution}

\begin{solution}{exo:b1:species-interactions:11}
जब फ़ॉस्फ़ेट वृद्धि सीमित करता हो, तब फ़ॉस्फ़ेट का \qty{80}{\%}
शर्करा के \qty{15}{\%} से कहीं अधिक मोल का है: यानी \omterm{def:b1:species-interactions:types}{सहोपकारिता}। किसी
फ़ॉस्फ़ेट-धनी मिट्टी में पौधा फ़ॉस्फ़ेट स्वयं ले लेता और कवक
कुछ न देकर \qty{15}{\%} लेता: यानी \omterm{def:b1:species-interactions:types}{परजीविता}। प्रयोग: फ़ॉस्फ़ेट के एक परास
पर कवक के साथ और बिना पौधे उगाइए और जैवभार नापिए; यह साहचर्य वहाँ
सहोपकारी है जहाँ कवकमूली पौधे बड़े उगें, और \omterm{def:b1:species-interactions:symbiosis}{परजीवी} वहाँ जहाँ वे छोटे उगें।
\end{solution}

\begin{solution}{exo:b1:species-interactions:12}
बाड़-भूखंड: नदियों के किनारे बाड़ वाले ऐसे भूखंड जिनमें एल्क घुस न सकें,
जिनके साथ खुले भूखंड जोड़े गए हों, जो भेड़ियों के झुंड वाली और बिना वाली
कई घाटियों में दोहराए गए हों, और जिन्हें फिर से बसाने से पहले तथा बाद
वर्षों तक नापा गया हो; विलो की ऊँचाई तथा आवरण, एल्क का घनत्व और हर भूखंड
में बिताया समय (पगचिह्न, कैमरे), शिकार और बर्फ़ की गहराई दर्ज कीजिए।
यह अनुप्रपात यह भविष्यवाणी करता है कि खुले भूखंडों में विलो केवल वहीं
उबरें जहाँ भेड़िए हों, और बाड़ के भीतर जितनी ही दर से; और यदि विलो
भेड़ियों के साथ और बिना बराबर उबरें, या बर्फ़ तथा शिकार का पीछा करें, तो
अनुप्रपात गलत ठहरता है या उन कारकों में बँट जाता है।
\end{solution}

\begin{solution}{pb:b1:species-interactions:1}
\textbf{1.} A: $N_A + 1.4 N_B = 200$, अंतःखंड 200 ($N_A$) और 143
($N_B$)। B: $N_B + 0.9 N_A = 130$, अंतःखंड 130 ($N_B$) और 144
($N_A$)।
\textbf{2.} अंतःखंड ($N_A$ अक्ष पर 144 के सामने 200, $N_B$ अक्ष पर
130 के सामने 143) A की \omterm{prop:b1:species-interactions:lotka}{सम-रेखा} को पूरी तरह B की \omterm{prop:b1:species-interactions:lotka}{सम-रेखा} से बाहर रख देते
हैं: $N_A = 200 - 1.4 N_B$ को $N_B = 130 - 0.9 N_A$ के साथ हल करने पर
$N_A = -69$ मिलता है, इसलिए रेखाएँ धनात्मक चतुर्थांश के बाहर मिलती हैं।
A किसी भी शुरुआती बिंदु से जीतती है: जहाँ भी B अब भी बढ़ सकती है वहाँ A
भी बढ़ती है और उससे अधिक टिकती है।
\textbf{3.} $T_A = 0.87$ दिन, $T_B = 1.26$ दिन: यानी A, जो तेज़
बढ़ती है।
\textbf{4.} A: $N_A + 0.3 N_C = 200$ (200, 667); C: $N_C + 0.3 N_A =
100$ (100, 333)। हर जाति स्वयं को दूसरी से अधिक सीमित करती है
($0.3 < 200/100$ और $0.3 < 100/200 = 0.5$): यानी स्थायी सहअस्तित्व।
\textbf{5.} कोई वन अचर नहीं है (ऋतुएँ, विक्षोभ दौड़ को तय होने से
पहले ही फिर से बैठा देते हैं) और उसमें बहुत-से संसाधन तथा जगहें हैं, इसलिए
जातियाँ एक ही को बाँटने के बजाय \omterm{def:b1:ecosystem-organization:niche}{निकेत} बाँट लेती हैं।
\textbf{6.} $N_A = 200 - 0.3 N_C$, $N_C = 100 - 0.3 N_A = 100 - 60 +
0.09 N_C$, इसलिए $N_C = 44$ और $N_A = 187$।
\textbf{7.} चढ़ती और लगभग 10 की ओर चपटी होती: यानी प्रकार II.
\textbf{8.} $(1/N, 1/f)$: $(0.2, 0.5)$, $(0.1, 0.303)$, $(0.05, 0.2)$,
$(0.025, 0.149)$, $(0.0125, 0.125)$: यानी एक सीधी रेखा।
\textbf{9.} ढाल $(0.5 - 0.125)/(0.2 - 0.0125) = 2.0 = 1/a$, इसलिए $a =
\qty{0.5}{m^2/d}$; अंतःखंड $0.5 - 2\times 0.2 = 0.1$ दिन $= h$।
\textbf{10.} अधिकतम $1/h = 10$ शिकार प्रति दिन; अर्ध-संतृप्ति $N =
1/(ah) = 20$ प्रति वर्ग मीटर पर।
\textbf{11.} निपटान अंश $= f h = f/10$: यानी आधे से अधिक तब जब $f
> 5$ हो, यानी प्रति वर्ग मीटर 40 और 80 पर (20 पर 5.0 ठीक आधा है)।
\textbf{12.} शिकार की वृद्धि $0.1 N (1 - N/200)$ खपत
$0.5N/(1 + 0.05N)$ के बराबर: $0.1(1 - N/200) = 0.5/(1 + 0.05 N)$, यानी
$(1 - N/200)(1 + 0.05N) = 5$: $1 + 0.05N - N/200 - N^2/4000 = 5$, इसलिए
$N^2 - 180 N + 16\,000 = 0$, $N = 90 \pm \sqrt{8100 - 16\,000}$: कोई
वास्तविक मूल नहीं --- प्रति वर्ग मीटर एक भृंग उससे अधिक खाता है जितना
शिकार कभी बना सकता है (अधिकतम वृद्धि $rK/4 = 5$ प्रति दिन, $N = 100$ पर,
जहाँ भृंग 8.3 खाता है) और उसे विलुप्ति तक ले जाता है। कोई साम्य नहीं;
शिकार तभी टिकता है जब भृंग कम हों या शिकार के पास शरणस्थल हों।
\textbf{13.} दूसरे शिकार के लिए \omterm{def:b1:species-interactions:functional}{निपटान काल} आधा हो जाता है ($h' = 0.05$),
इसलिए संतृप्ति पर भृंग दिन में उनमें से 20 तक खा लेता है: ऊँचे घनत्व पर
कुल ग्रहण मोटे तौर पर दुगना हो जाता है यदि वह तेज़ निपटने वाले शिकार की
ओर मुड़ जाए।
\textbf{14.} $H_{15} = \ln 15 = 2.71$। आठ जातियाँ: $-(0.85\ln 0.85
+ 7\times 0.0214\ln 0.0214) = 0.138 + 0.576 = 0.71$।
\textbf{15.} $2.71/0.71 = 3.8$ के गुणक से।
\textbf{16.} सीपी पर तारामीन के $-$ ने सीपी के अपने स्पर्धियों पर $-$ को
कमज़ोर रखा; पहला हटाइए और सीपी उस घनत्व तक पहुँच जाती है जिस पर उसकी
\omterm{def:b1:species-interactions:types}{स्पर्धा} बाकियों को अपवर्जित कर देती है।
\textbf{17.} जैवभार यह नापता है कि कोई जाति ऊर्जा-बहाव का कितना रखती है,
यह नहीं कि उसकी अन्योन्यक्रियाएँ क्या करती हैं; किसी \omterm{def:b1:species-interactions:keystone}{आधार जाति} का असर
उन जातियों से होकर जाता है जिन्हें वह काबू में रखती है, और जिनका जैवभार
बड़ा है।
\textbf{18.} तारामीन हटाने से कोई ऐसी दुर्लभ जाति छूटती जो चट्टान पर
कब्ज़ा न कर पाती: सीपी-प्रधान \omterm{def:b1:ecosystem-organization:ecosystem}{समुदाय} में कम ही बदलाव होता, और तारामीन
कोई \omterm{def:b1:species-interactions:keystone}{आधार जाति} न होता।
\textbf{19.} छोटे धब्बों से तारामीन को बाहर रखते पिंजरे, और पिंजरे के
अपने असर के नियंत्रण के लिए खुले पिंजरे: यदि सीपियाँ केवल बंद पिंजरों में
कब्ज़ा करें, तो सीपियों पर \omterm{def:b1:species-interactions:types}{परभक्षण} ही क्रियाविधि है; और ऐसे पिंजरे
जिनमें तारामीन हो पर सीपियाँ हाथ से हटा दी गई हों यह जाँचेंगे कि तारामीन
के दूसरे शिकार मायने रखते हैं या नहीं।
\textbf{20.} कमाई $\qty{40}{mg}$, लागत $8\times 0.5 = \qty{4}{mg}$: शुद्ध
\qty{36}{mg} $=$ \qty{0.58}{kJ}।
\textbf{21.} \qty{4}{mg} में से \qty{40}{\micro g}: यानी फूल के दैनिक
प्रकाश-संश्लेषित पदार्थ का \qty{1}{\%} लगभग तीन पराग-स्थानांतरण, यानी
पौधे का जनन खरीदता है।
\textbf{22.} हर एक दूसरे से कुछ लेता है, पर हर एक जितना खोता है उससे
अधिक कमाता है: चिह्न शुद्ध प्रभाव का होता है, दोहन का
नहीं।
\textbf{23.} चोर $+$, पौधा $-$ (मकरंद गया, कोई परागण नहीं): यानी उस
\omterm{def:b1:species-interactions:types}{सहोपकारिता} की एक \omterm{def:b1:species-interactions:types}{परजीविता}; और मधुमक्खी को खाली फूल मिलते हैं तथा वह कम
कमाती है, यानी एक परोक्ष $-$।
\textbf{24.} पौधा $\leftrightarrow$ मधुमक्खी $+/+$; चोर $\to$ पौधा $-$,
चोर $\to$ मधुमक्खी $-$; पक्षी $\to$ चोर $-$। दो ऋण चिह्न: पक्षी का
पौधे पर और मधुमक्खी पर परोक्ष $+$ है।
\textbf{25.} $a = \qty{0.5}{m^2/d}$, $h = \qty{0.1}{d}$ (\qty{2.4}{h}),
और अधिकतम प्रति दिन \num{10} शिकार।
\end{solution}
